\chapter*{Resumen}

La proliferación de los dispositivos de geolocalización ha transformado el estudio
del movimiento animal: el cuello de botella de la disciplina se ha desplazado de
registrar dónde ha estado un animal a anticipar dónde estará. Este trabajo aborda
esa pregunta sobre una especie concreta, la gaviota sombría (\textit{Larus fuscus}),
un migrante de larga distancia cuyo corredor une las zonas de cría del norte de
Europa con las áreas de invernada del este de África, a partir de datos públicos de
seguimiento GPS. El problema se plantea a escala diaria: dada la posición de un
individuo hoy, predecir dónde estará mañana.

Para ello se diseña y evalúa un enfoque híbrido que combina modelos probabilísticos
y aprendizaje automático. A partir de los registros GPS en bruto se construye una
secuencia diaria homogénea de posiciones por individuo, con los huecos del
seguimiento representados de forma explícita, que sirve de base común a todos los
modelos. Sobre ella se articulan tres piezas. Una cadena de Markov visible establece
una baseline interpretable y fija el criterio de evaluación común, frente a la
referencia trivial de suponer que el animal permanecerá donde está. Un modelo oculto
de Markov infiere, sin supervisión, el régimen de comportamiento de cada día
(estacionario o en migración) y se valida por su coherencia con la fenología de la
especie. Ese régimen alimenta, como una variable más, a modelos de aprendizaje
supervisado (Random Forest, XGBoost y LightGBM), que predicen el desplazamiento
empleando únicamente información disponible antes del instante que se predice. Por
último, una herramienta de cartografía interactiva permite visualizar y comparar
sobre el mapa las predicciones de los distintos modelos.

El enfoque aporta en tres frentes complementarios. Calibra la incertidumbre de cada
predicción mediante una banda que la herramienta interactiva traslada al mapa; aísla
el régimen de migración, que es donde anticipar el desplazamiento resulta
verdaderamente informativo, frente a los largos periodos estacionarios propios de la
especie; y caracteriza con datos el techo estructural que impone predecir a un día con
información puramente local, un límite inherente al horizonte del problema más que a un
modelo concreto. La formulación del problema como regresión de cuantiles es la que
mejor concilia la precisión de la posición estimada con una incertidumbre bien
calibrada.

La contribución del trabajo es triple: la caracterización rigurosa de un límite
estructural de la predicción a corto plazo, una herramienta interactiva para comparar
modelos sobre el territorio, y un criterio metodológico transversal, el de exigir
coherencia biológica a todo resultado antes de darlo por válido. De ahí se derivan las
líneas de trabajo futuro, que apuntan a modelos de secuencia, a la incorporación del
viento y a la predicción a varios pasos.

\begin{otherlanguage}{english}
  \chapter*{Abstract}

  \todo[inline]{Abstract of the Final Degree Project. Maximum length: 2 pages.}
\end{otherlanguage}


%%%%%%%%%%%%%%%%%%%%%%%%%%%%%%%%%%%%%%%%%%%%%%%%%%%%%%%%%%%
%% Final del resumen.
%%%%%%%%%%%%%%%%%%%%%%%%%%%%%%%%%%%%%%%%%%%%%%%%%%%%%%%%%%%
